\textbf{COLUNI 2026 -- Questão 32}\par\medskip Leia a notícia a seguir.

“Trump assina decreto que impõe tarifa de 50\% ao Brasil, mas com várias exceções. Segundo a Casa Branca, a medida é uma resposta a ações do governo brasileiro que representariam uma ‘ameaça incomum e extraordinária à segurança nacional’. Aeronaves civis, suco de laranja e petróleo estão entre os itens que ficaram de fora.”\par
(Disponível em: \url{https://g1.globo.com/economia/noticia/2025/07/30/trump-assina-decreto-que-impoe-tarifa-de-50percent-ao-brasil.ghtml}. Adaptado. Acesso em: 30 jul. 2025.)

Sobre essa decisão e seus possíveis significados, analise as afirmações abaixo:
\begin{enumerate}
\item[I.] A decisão dos EUA mostra como os países usam o comércio internacional para pressionar politicamente outros governos.
\item[II.] A retirada de alguns produtos da tarifa pode indicar que os Estados Unidos têm interesse em continuar comprando ou vendendo esses itens.
\item[III.] Tarifas como essas aplicadas pelos EUA tornam os produtos brasileiros mais baratos para os consumidores estadunidenses, incentivando mais importações.
\item[IV.] A medida de Trump pode causar conflitos econômicos e prejudicar a relação entre os dois países.
\end{enumerate}
Está CORRETO o que se afirma em:\begin{enumerate}
\item I, II e IV, apenas.
\item I e III, apenas.
\item II e IV, apenas.
\item I, II, III e IV.\end{enumerate}